\section{Conclusion}
\label{sec:conclusion}

% NOTE (2026-09-25): this sentence previously opened with "a parameter- and
% memory-efficient method". The memory half is not merely unsupported but BACKWARDS:
% FOLoRA's importance model is a rank-$k$ kernel and is 32.6x EWC's per-layer memory at
% k=64 (182.8 MB vs 5.6 MB; see 内部决策日志 11.1), which the same conclusion
% correctly reports three sentences later. An opening that claims the opposite of the
% paper's own measurement is the first thing a referee sees. Do not reinstate
% "memory-efficient"; the only efficiency claim this paper earns is that the
% *trainable-parameter* count does not grow with the number of tasks.
We propose FOLoRA, a replay-free method for continual learning that mitigates catastrophic
forgetting in low-rank adaptation while keeping the number of trainable parameters fixed
as the number of tasks grows. FOLoRA maintains an accumulated
second-order importance kernel for each adapted layer, forms an importance-weighted
protected subspace from its top eigendirections, and penalizes alignment of the adapter's
displacement with those directions in proportion to their importance. We derive a
forgetting upper bound that decomposes into Fisher-weighted subspace overlaps and showed
that FOLoRA's regularizer minimizes its dominant term, thereby providing a principled
stability--plasticity trade-off and containing equal-weight orthogonalization as a special
case. We also measure the slack of that bound instead of assuming it away: at the
operating point we probe the remainder and the stationarity residual are the same order as
the quadratic term, so the bound is presented as a design principle for what to penalize
rather than as a tight or numerically predictive characterization of forgetting. The
empirical case for the regularizer rests on the $\lambda{=}0$ ablation, which removes the
Fisher-weighted penalty while changing nothing else and returns the method to the level of
unprotected sequential fine-tuning. Experiments on standard class-incremental benchmarks show that FOLoRA reaches an
average accuracy on par with the strongest regularized baseline while using no replay
buffer and no growth in the parameter count, and we quantify the importance-model memory
it pays for representing that model as a rank-$k$ subspace rather than EWC's diagonal one
--- $32\times$ EWC's importance state at $k{=}64$ and $8.5\times$ at $k{=}16$, which is the
trade-off that makes $k$ a memory--accuracy dial rather than a free parameter. We also
report the run-to-run reproducibility floor we measured ($1.29$ points at a fixed seed and
a fixed configuration), and we state the resulting conclusions at the strength the data
supports: parity with the strongest regularized baseline, rather than an improvement over
it, and a significant advantage over the parameter-efficient prompt and subspace
baselines.
% NOTE (2026-09-24): the sentence here previously read "...demonstrate CONSISTENT
% IMPROVEMENTS over sequential LoRA, EWC-LoRA, O-LoRA, L2P and CODA-Prompt in both average
% accuracy and forgetting". That is FALSE under the NCM protocol the paper reports: at the
% selected operating point FOLoRA is at parity with EWC-LoRA (paired t = 0.94, p = 0.41 at
% n=4; see 内部决策日志 13.3), and at other points in the lambda grid it is BELOW it
% (lambda=300, k=16: 73.61 vs 77.22). It also contradicted the abstract, which had already
% been corrected to "on par with the strongest regularized baseline". A conclusion is where
% referees check whether the paper claims more than its tables show -- do NOT restore the
% "improvements" wording, and do not add a forgetting-reduction claim either: under NCM,
% ACC and FGT co-vary across methods, so a forgetting reduction is not an independently
% earned result (same red line as 05_experiments.tex).

A natural direction for future work is to extend FOLoRA to task-agnostic settings with
distribution shift (domain-incremental learning) and to language models, where the
low-rank subspace geometry of different task families can be exploited more explicitly.
